% Homework Bridge --- one-pager
% Suvidha International Foundation AI Virtual Hackathon 2026
% Build with:  pdflatex ONEPAGER.tex
\documentclass[9pt,letterpaper]{extarticle}

\usepackage[margin=0.42in,top=0.32in,bottom=0.30in]{geometry}
\usepackage[T1]{fontenc}
\usepackage[utf8]{inputenc}
\usepackage{charter}
\usepackage{xcolor}
\usepackage{multicol}
\usepackage{enumitem}
\usepackage{titlesec}
\usepackage{parskip}
\usepackage{microtype}
\usepackage[hidelinks]{hyperref}

\definecolor{ink}{HTML}{1C1917}
\definecolor{accent}{HTML}{9A4A2F}
\definecolor{muted}{HTML}{6B6459}
\definecolor{rule}{HTML}{D9D2C5}
\definecolor{band}{HTML}{F4EBE5}

\color{ink}
\pagestyle{empty}
\setlength{\columnsep}{15pt}

\titleformat{\section}{\bfseries\large\color{accent}}{}{0pt}{}
\titlespacing*{\section}{0pt}{4pt}{1pt}

\begin{document}

%% ---------- masthead ----------
{\fontsize{23}{25}\selectfont\bfseries Homework Bridge\par}
\vspace{1pt}
{\color{accent}\rule{\textwidth}{1.2pt}}
\vspace{4pt}

{\fontsize{13}{15}\selectfont\bfseries\itshape My dad always knew how to do my math homework. He just couldn't get it to me.\par}
\vspace{5pt}

\noindent\colorbox{accent}{\parbox{\dimexpr\textwidth-2\fboxsep}{\vspace{4pt}\large\color{white}%
\textbf{AI disclosure:} Claude (Anthropic) wrote the code and helped polish this write-up.
I directed the project --- the idea, the research into how other countries teach math, the
design decisions, and the conversation with my dad are mine.\vspace{4pt}}}
\vspace{3pt}

\noindent\colorbox{band}{\parbox{\dimexpr\textwidth-2\fboxsep}{\vspace{3pt}%
\textbf{Live demo:} \url{https://homework-bridge-rho.vercel.app} \hfill
\textbf{Code:} \url{https://github.com/Adarsh-Samanu/homework-bridge}\vspace{3pt}}}
\vspace{4pt}

\begin{multicols}{2}

\section*{The problem}

My dad grew up in India and learned all of this math there --- long division, subtraction with borrowing. He is good at it. He speaks Telugu, and English is hard for him. I never learned enough Telugu for him to just explain it to me instead.

So when I got stuck on a worksheet in elementary and middle school, this is what it looked like: I'd be sitting there with the paper, my dad right next to me, and he \emph{knew the answer} --- and there was no way for it to get from him to me. He couldn't read the sheet. I couldn't follow his explanation. One language apart, at the same table.

I asked him about it while building this. He said something like this would have really helped him back then --- he wanted to teach me the way he had learned it, and couldn't.

\textbf{And translating the worksheet would not have fixed us.} Hand him my homework in perfect Telugu and he still can't help, because it says things like \emph{``use a number bond to decompose the second number.''} There is no Telugu word for \textbf{number bond} --- not because it's a hard translation, but because the thing doesn't exist in what he learned. He was taught column subtraction with borrowing. He'd be staring at a diagram he had never seen, now in Telugu.

It isn't only vocabulary. \textbf{Long division sits somewhere else on the page} in Mexico, Vietnam, and Brazil. \textbf{\texttt{3,5} means three and a half} in most of Latin America, so a parent can read \texttt{3.5} as thirty-five and never find out. \textbf{\texttt{1,00,000} is one lakh} in India. \textbf{\texttt{187 r 1}} --- nothing says the \texttt{r} means remainder.

TalkingPoints, ParentSquare, ClassDojo and Google Lens all translate school communication well. Every one would have left my dad exactly where he was.

\section*{The solution}

Photograph the worksheet or paste it in, pick your language and the country you went to school in. Then, in your language:

\begin{enumerate}[leftmargin=13pt,itemsep=1pt,topsep=2pt]
  \item \textbf{What the homework is actually asking} --- plainly, not word-for-word
  \item \textbf{The way the school wants it done}, step by step
  \item \textbf{The same problem done the way \emph{you} were taught} $\leftarrow$ nothing else does this
  \item \textbf{How the two line up}, and the one place they really differ
  \item \textbf{Notation warnings} --- the comma and decimal traps above
  \item \textbf{Three questions to ask your kid} that check understanding without giving the answer
\end{enumerate}

Worked examples always use \textbf{different numbers} than the real homework, so you can't copy them onto the sheet. This helps a parent help; it doesn't do the homework.

\vspace{2pt}
{\color{accent}\textbf{Try the third sample --- Telugu, India, number bonds.} That is my dad's exact situation.}

\section*{Tech stack}

\textbf{Front end \& hosting.} Next.js 16 (App Router), React 19, TypeScript, deployed on Vercel. Read-aloud uses the browser's Web Speech API --- free and offline. No database, no accounts, nothing stored.

\textbf{Models}, through Featherless AI's OpenAI-compatible API: \texttt{GLM-5.2} does all reasoning and writes the explanation in the parent's language; \texttt{Qwen3-VL} transcribes photos and is never allowed to solve anything. An Anthropic-SDK adapter is swappable via one environment variable. Every stage walks a fallback chain, because individual models go busy without warning.

\textbf{The part that matters most is a data file.} Asked directly, the model gave vague and inconsistent answers about how other countries teach arithmetic. So we researched the real differences and wrote them down country by country in \texttt{lib/methods.ts}. Deciding that file had to exist is the call that makes the comparison trustworthy.

\textbf{Correctness is machine-checked.} The first working version told a Telugu-speaking parent that $62-27=55$ --- confidently, in the panel a parent is meant to trust. A homework helper that gets numbers wrong is worse than none, because now the parent looks wrong in front of their kid. So reading and reasoning were split apart, and \texttt{eval/run.py} now re-computes every equation the model writes and fails the run on any mismatch. The three demo worksheets pass: 8, 16 and 8 equations checked, zero wrong.

\section*{What's next}

\begin{itemize}[leftmargin=11pt,itemsep=1pt,topsep=2pt]
  \item \textbf{More countries.} Six are covered. The Philippines, Haiti, Nigeria and Central America are the obvious gaps.
  \item \textbf{Handwriting.} Printed worksheets read reliably; handwritten ones are the weakest part.
  \item \textbf{Test with real families,} not just mine. One interview shaped this; ten would shape it better.
  \item \textbf{Speed.} Answers take 8--20 seconds live. The demo worksheets are pre-rendered because of it.
  \item \textbf{Beyond math} --- science and reading have their own untranslatable methods.
\end{itemize}

\end{multicols}

\vspace{2pt}
{\color{rule}\rule{\textwidth}{0.5pt}}
\vspace{1pt}

\noindent{\small\color{muted}\textbf{Demo:} homework-bridge-rho.vercel.app \quad\textbullet\quad \textbf{Repo:} github.com/Adarsh-Samanu/homework-bridge \hfill Suvidha AI Virtual Hackathon 2026}

\end{document}
